\section*{Capítulo \ref{ch:b1:stratigraphic-scale} --- A escala estratigráfica}

\begin{solution}{exo:b1:stratigraphic-scale:1}
(a)~Cronoestratigráfica, um \omterm{def:b1:stratigraphic-scale:chronostratigraphy}{sistema}. (b)~Litoestratigráfica, uma \omterm{def:b1:stratigraphic-scale:lithostratigraphy}{formação}.
(c)~Geocronológica, uma época (sua rocha é a \omterm{def:b1:stratigraphic-scale:chronostratigraphy}{Série} do Cretáceo Inferior).
(d)~Bioestratigráfica, uma \omterm{def:b1:stratigraphic-scale:biozone}{biozona}. (e)~Cronoestratigráfica, um \omterm{def:b1:stratigraphic-scale:chronostratigraphy}{andar}.
(f)~Geocronológica, uma época. (g)~Litoestratigráfica, um \omterm{def:b1:stratigraphic-scale:lithostratigraphy}{membro}.
\end{solution}

\begin{solution}{exo:b1:stratigraphic-scale:2}
Paleozoico $538.8 - 251.902 = \qty{286.9}{\Myr}$, \qty{53.2}{\%} do
Fanerozoico; Mesozoico $251.902 - 66.00 = \qty{185.9}{\Myr}$, \qty{34.5}{\%};
Cenozoico \qty{66.0}{\Myr}, \qty{12.2}{\%}. O Fanerozoico é
$538.8/4567 = \qty{11.8}{\%}$ do tempo desde a base do Hadeano.
\end{solution}

\begin{solution}{exo:b1:stratigraphic-scale:3}
A, B, C e D foram depositadas nessa ordem (superposição), horizontalmente (horizontalidade
original), de modo que sua inclinação é um evento posterior. O dique E corta A a D, então é
mais jovem que D (relações de corte); se veio antes ou depois da inclinação não se pode dizer
pela descrição. A superfície de erosão corta D e E, então é mais jovem que ambos (relações de
corte). F e G foram então depositadas sobre ela (superposição). A falha desloca A a F, mas não
G: moveu-se depois de F e antes de G (relações de corte). Ordem: A, B, C, D; inclinação e
intrusão de E (ordem indecisa); erosão; F; falhamento; G.
\end{solution}

\begin{solution}{exo:b1:stratigraphic-scale:4}
$r = 30/(85.40 - 84.60) = \qty{37.5}{m}$ por milhão de anos. A primeira aparição
fica $12/37.5 = \qty{0.32}{\Myr}$ depois da cinza inferior: $85.40 - 0.32 =
\qty{85.08}{\Ma}$.
\end{solution}

\begin{solution}{exo:b1:stratigraphic-scale:5}
(a)~Quatro barras verticais em função da altura. (b)~B e C ocorrem juntos de
\qtyrange{26}{50}{m}: uma \omterm{def:b1:stratigraphic-scale:biozone}{zona de distribuição} concorrente de \qty{24}{m} de espessura. A \omterm{def:b1:stratigraphic-scale:biozone}{zona
de intervalo} da primeira aparição de B à de D vai de \qtyrange{12}{41}{m}, \qty{29}{m}. (c)~Na
segunda seção as primeiras aparições de B (\qty{3}{m}) e de C (\qty{9}{m}) podem ser situadas,
de modo que a base da \omterm{def:b1:stratigraphic-scale:biozone}{zona de intervalo} de (b) e a base da \omterm{def:b1:stratigraphic-scale:biozone}{zona de distribuição} concorrente são
reconhecidas; seu topo (a primeira aparição de D) e o topo da sobreposição como visto na
seção~1 não podem ser verificados, pois D está ausente. As primeiras aparições de B e C
($\qty{12}{m} \leftrightarrow
\qty{3}{m}$, $\qty{26}{m} \leftrightarrow \qty{9}{m}$) são as linhas mais seguras; as últimas
aparições (B a \qty{50}{m} e \qty{15}{m}) dependem mais da preservação.
\end{solution}

\begin{solution}{exo:b1:stratigraphic-scale:6}
(a)~$60/20 = \qty{3}{\Myr}$: o topo do arenito é três milhões de anos mais jovem no poço do
continente. (b)~Um milhão de anos de lama a \qty{30}{m} por milhão de anos: a cinza fica
\qty{30}{m} acima do topo do arenito. (c)~Quando a cinza caiu, a linha de costa tinha
avançado \qty{20}{km} para o interior a partir do poço do mar, de modo que ainda estava
\qty{40}{km} mar adentro do poço do continente, que estava em terra seca. A cinza caiu sobre a
superfície da terra ali, abaixo do arenito que o mar depositou dois milhões de anos depois ---
se não tiver sido levada antes pela erosão, como uma camada solta em terra pode muito bem ser.
A mesma camada de cinza está acima do limite de \omterm{def:b1:stratigraphic-scale:lithostratigraphy}{formação} em um poço e abaixo dele no outro: o
limite atravessa o tempo.
\end{solution}

\begin{solution}{exo:b1:stratigraphic-scale:7}
(a)~Lochkoviano $1.36/419.62 = \qty{0.32}{\%}$; Induano $0.024/251.902 =
\qty{0.0095}{\%}$, cerca de 34 vezes menor. A base do Induano fica entre camadas de cinza
datadas com precisão na própria seção do GSSP; uma incerteza grande sugere um limite cuja idade
deve ser transportada de níveis datados mais distantes. (b)~$251.941 - 251.880 = \qty{0.061}{\Myr} = 61$ mil anos, com a incerteza $\sqrt{37^2 + 31^2} = 48$ mil anos: os publicados
\num[uncertainty-mode=separate]{60 \pm 48} mil anos.
\end{solution}

\begin{solution}{exo:b1:stratigraphic-scale:8}
(a)~$m = 3000 \times \frac43\pi\,(5000)^3 = \qty{1.57e15}{kg}$; a poeira,
$60m = \qty{9.42e16}{kg}$. (b)~A superfície da Terra é $4\pi(6.378\times10^6)^2
= \qty{5.11e14}{m^2}$: \qty{184}{kg} por metro quadrado. (c)~$184/2500 =
\qty{0.074}{m}$, cerca de \qty{7}{cm}. (d)~A espessura varia como o cubo do
diâmetro: \qty{1.6}{cm} para \qty{6}{km}, \qty{20}{cm} para \qty{14}{km}. A ordem
de grandeza, de centímetros a decímetros, é a da argila de fronteira, de \qty{50}{cm} de
espessura em El Kef; a comparação só pode ser grosseira, já que a poeira não pode ter caído de
modo uniforme, e uma argila também reúne o sedimento comum de seu tempo.
\end{solution}

\begin{solution}{exo:b1:stratigraphic-scale:9}
(a)~$\num{60e3} \times \num{2e-6} = \qty{0.12}{m}$, \qty{12}{cm}; para 12 e 108 mil anos,
\qty{2.4}{cm} e \qty{21.6}{cm}. (b)~$\num{32e3} \times
\num{10e-6} = \qty{0.32}{m}$. (c)~Uma extinção ocupa de centímetros a decímetros de rocha: as
amostras devem ser tomadas centímetro a centímetro, e uma única amostra de \qty{10}{cm} na seção
condensada representa, em média, cinquenta mil anos --- a maior parte do evento.
\end{solution}

\begin{solution}{exo:b1:stratigraphic-scale:10}
(a)~$u = R/s = 10 \times 2000 = \qty{20000}{m}$ por milhão de anos,
\qty{20}{km} por milhão de anos; em \qty{30}{km}, $\Delta t = \qty{1.5}{\Myr}$.
(b)~A planície é agora a superfície $z = s x + a t$. A linha de costa satisfaz
$s x + a t = z_0 + R t$, logo $x(t) = (z_0 + (R - a)t)/s$ e
$\mathrm dx/\mathrm dt = (R - a)/s$, positiva (para o interior) enquanto $a < R$. Quando
$a > R$ a velocidade é negativa: a linha de costa se move para o mar, uma regressão, e o
arenito da linha de costa é então mais jovem em direção ao mar. (c)~Uma falésia de
\qty{200}{m} de comprimento abrange, a \qty{20}{km} por milhão de anos, $0.2/20 = \qty{0.01}{\Myr}$
--- dez mil anos de diacronismo, muito abaixo do que fósseis ou a maioria das datas podem
resolver; as rochas dos dois lados do contato são as mesmas em toda a sua extensão.
\end{solution}

\begin{solution}{exo:b1:stratigraphic-scale:11}
(a)~$f = (25 - 10)/(50 - 10) = 0.375$ e $t = 100.50 - 0.375 \times 0.80 =
\qty{100.20}{\Ma}$, com
\[
  \sigma_t = \sqrt{(0.625 \times 0.05)^2 + (0.375 \times 0.03)^2} = \qty{0.033}{\Myr}.
\] (b)~$g(f) = (1 - f)^2\sigma_1^2 + f^2\sigma_2^2$
tem $g'(f) = -2(1 - f)\sigma_1^2 + 2f\sigma_2^2$, zero em $f^* =
\sigma_1^2/(\sigma_1^2 + \sigma_2^2) = 0.0025/0.0034 = 0.735$, mais perto da melhor
camada datada; ali $\sigma_t = \sigma_1\sigma_2/\sqrt{\sigma_1^2 + \sigma_2^2} =
\qty{0.026}{\Myr}$, menor que ambas. (c)~$g$ é uma parábola voltada para cima, então
em $[0, 1]$ ela é maior em uma extremidade: $g(0) = \sigma_1^2$, $g(1) = \sigma_2^2$.
Logo $\sigma_t \le \max(\sigma_1, \sigma_2)$ entre as camadas.
\end{solution}

\begin{solution}{exo:b1:stratigraphic-scale:12}
(a)~O limite está na metade da profundidade: $f = 0.5$, \qty{29.50}{\Ma}. (b)~Com
$\varphi_0\lambda = 0.63 \times 1960 = \qty{1234.8}{m}$, a espessura sólida de
\qtyrange{600}{1000}{m} é $400 - 1234.8\,(e^{-600/1960} - e^{-1000/1960}) =
\qty{232.2}{m}$, e de \qtyrange{200}{600}{m}, \qty{194.2}{m}; ao todo
\qty{426.3}{m}. A fração dos grãos depositados entre a cinza inferior e o limite é
$232.2/426.3 = 0.545$, e a idade é $30.0 - 0.545 =
\qty{29.46}{\Ma}$. (c)~A segunda estimativa é \num{45} mil anos mais jovem. É a melhor: a
metade mais profunda perdeu mais água, de modo que contém mais grãos, e mais tempo, por metro
do que a metade mais rasa (porosidade média \num{0.42} contra \num{0.51}).
\end{solution}

\subsection*{Problema: correlacionando três seções através de uma bacia}

\begin{solution}{pb:b1:stratigraphic-scale:1}
\begin{enumerate}
  \item Três colunas em, digamos, \qty{1}{cm} para \qty{10}{m}, unidas pelas camadas de cinza e
        pelos limites de zona; a cunha do Arenito Kell afina de \qty{45}{m} em A para
        \qty{6}{m} em C.
  \item Zona~2: A $65.8 - 36.6 = \qty{29.2}{m}$, B \qty{20.8}{m}, C
        \qty{10.0}{m}. Zona~3 descrita: A \qty{14.2}{m}, B \qty{14.0}{m}, C
        \qty{7.2}{m}.
  \item A zona~1 começa na primeira aparição de \emph{V.~primum}, que ocorre desde a base
        de todas as seções: sua primeira aparição fica abaixo delas.
  \item Ele vive apenas no arenito: segue a fácies, e a fácies tem idades diferentes de lugar
        para lugar. Em A ocorre nas zonas~1 e~2, em B e C apenas na zona~1.
  \item Na zona~2 nas três (A \qty{49.0}{m}, B \qty{40.0}{m}, C \qty{19.0}{m}). É consistente
        de seção para seção aqui, e assim poderia subdividir a zona~2, mas uma última aparição
        depende da preservação; deveria ser verificado em mais seções antes.
  \item As linhas unem \qty{36.6}{m}, \qty{31.2}{m} e \qty{14.8}{m}, depois \qty{65.8}{m},
        \qty{52.0}{m} e \qty{24.8}{m}. A zona~2 afina na razão $29.2 : 20.8 : 10.0$, cerca de
        $2.9 : 2.1 : 1$: se a zona durou o mesmo tempo em toda parte, o sedimento se acumulou
        cerca de três vezes mais depressa em A do que em C.
  \item A: zona~2 ($45.0 > 36.6$); B: zona~1 ($20.0 < 31.2$); C: zona~1
        ($6.0 < 14.8$).
  \item Não: ela está na zona~2 em A, na zona~1 em B e C, logo é mais jovem a oeste. O folhelho
        marinho se espalhou para oeste sobre a areia da linha de costa: a linha de costa se moveu
        para oeste, para o interior, em uma transgressão.
  \item Uma cinza cai sobre toda a bacia em um instante geológico e se assenta sobre qualquer
        sedimento; o topo do arenito segue a linha de costa em movimento e é \omterm{def:b1:stratigraphic-scale:diachronous}{diácrono}.
  \item Continuidade lateral (a camada continua através da bacia), e superposição (cinza~1
        abaixo da cinza~2 em todas as seções, o que verifica a \omterm{def:b1:cross-sections:log}{correlação}).
  \item A $(71.0 - 5.0)/1.20 = \qty{55.0}{m}$ por milhão de anos; B
        $48.0/1.20 = \qty{40.0}{}$; C $24.0/1.20 = \qty{20.0}{}$.
  \item A está mais perto da terra, onde o sedimento era suprido; menos dele chegou às águas
        mais profundas a leste.
  \item $0.1/55 = \num{1.8}$ mil anos em A, $\num{2.5}$ em B, $\num{5.0}$ em C: A resolve o
        tempo mais finamente.
  \item A $152.40 - 40.0/55.0 = \qty{151.67}{\Ma}$; B $152.40 - 12.0/40.0 =
        \qty{152.10}{\Ma}$; C $152.40 - 3.0/20.0 = \qty{152.25}{\Ma}$.
  \item C para B: $20/(152.25 - 152.10) = \qty{133}{km}$ por milhão de anos; B para A:
        $25/(152.10 - 151.67) = \qty{58.5}{km}$ por milhão de anos. Ela desacelerou.
  \item A $152.40 - 60.8/55.0 = \qty{151.29}{\Ma}$; B $152.40 - 44.0/40.0 =
        \qty{151.30}{\Ma}$; C $152.40 - 21.8/20.0 = \qty{151.31}{\Ma}$.
  \item Elas abrangem \num{16} mil anos, menos que a incerteza de cada uma (cerca de \num{28}
        mil anos, questão~19): a primeira aparição é síncrona em toda a bacia dentro da
        resolução das datas, o que a torna um bom índice.
  \item A $152.40 - 31.6/55.0 = \qty{151.826}{\Ma}$; B $152.40 - 23.2/40.0 =
        \qty{151.820}{\Ma}$; C $152.40 - 11.8/20.0 = \qty{151.810}{\Ma}$.
  \item $f = 0.479$, $0.483$, $0.492$; $\sigma_t = \sqrt{(1 - f)^2 + f^2} \times
        0.04 = \qty{0.028}{\Myr}$ em cada, \num{28} mil anos.
  \item Média \qty{151.819}{\Ma}; dispersão $151.826 - 151.810 = \num{16}$ mil anos, menor que
        a incerteza de cada idade: as três concordam.
  \item A interpolação supõe uma só taxa da cinza~1 à cinza~2. Em A esse intervalo contém
        arenito de linha de costa até \qty{45}{m} e folhelho acima; areia e lama não precisam ter
        se acumulado à mesma taxa, e se compactam de modo diferente
        (\cref{ch:b1:diagenesis-basins}). Se a areia se acumulou mais depressa, a parte arenosa
        contém menos tempo do que sua espessura sugere, e a primeira aparição é mais antiga do que
        a calculada --- a idade de A é de fato a mais antiga das três.
  \item B: a primeira aparição fica em folhelho marinho, \qty{11}{m} acima do topo do arenito,
        sem mudança de fácies perto do nível, e a taxa dá \num{2.5} mil anos por \qty{10}{cm}. A
        resolve o tempo melhor, mas coloca o nível no arenito da linha de costa, \qty{8.4}{m}
        abaixo de uma mudança de fácies; C é condensada, \num{5} mil anos por \qty{10}{cm}. C é a
        pior para a resolução, A para as fácies.
  \item A partir das idades médias das duas primeiras aparições, $151.819 - 151.302 =
        \qty{0.52}{\Myr}$: cerca de meio milhão de anos.
  \item O limite é \qty[uncertainty-mode=separate]{151.82(3)}{\Ma}.
\end{enumerate}
\end{solution}
